\chapter{SCADA: historian, HMI and alarms}
\label{ch:scada}

Supervisory control and data acquisition (SCADA) sits above the PLC. It does
not close loops; it records, displays and lets people intervene. In pyplant
the SCADA layer is three pieces: the historian collector, its query API and
the operator HMI with an ISA-18.2 alarm manager.

\section{The historian}

\subsection{Data flow}

\begin{figure}[htbp]
\centering
\scalebox{0.8}{%
\begin{tikzpicture}[node distance=7mm and 8mm]
  \node[hot] (opc) {PLC OPC UA\\server, 5403};
  \node[box,right=of opc] (sub) {subscription\\publish 200\,ms};
  \node[box,right=of sub] (buf) {buffer\\\code{_on_change}};
  \node[box,right=of buf] (db) {DuckDB\\\code{samples}};
  \node[box,below=of buf] (hb) {heartbeat loop\\re-store quiet tags\\every 30\,s};
  \node[good,below=of db] (api) {query API\\port 5404};
  \node[box,right=of api] (cli) {HMI trends,\\inference};
  \draw[flow] (opc) -- (sub);
  \draw[flow] (sub) -- node[lbl,above]{on change} (buf);
  \draw[flow] (buf) -- node[lbl,above]{flush 1\,s} (db);
  \draw[flow] (hb) -- (buf);
  \draw[flow] (db) -- (api);
  \draw[flow] (api) -- (cli);
\end{tikzpicture}}
\caption{The historian collector (\code{pyplant/services/historian.py}).}
\label{fig:historian}
\end{figure}

The collector subscribes to every tag the PLC serves. An OPC UA subscription
already pushes only on change, which is most of what a historian's
compression does. What it does not give is an anchor for a tag that has been
quiet for an hour, so each tag is also re-stored if it has been silent for
\code{heartbeat} seconds (default 30). This is the classic exception-reporting
pair: store on change, and never let the gap exceed a maximum. Writes are
buffered and flushed in batches every second, because row-at-a-time inserts
are the usual way to make DuckDB look slow.

\subsection{Storage}
The table is in long format, one row per sample:
\begin{center}
\code{samples(ts TIMESTAMP, tag VARCHAR, value DOUBLE, quality UTINYINT)}
\end{center}
A wide table needs a migration the day a tag is added; a long table does not.
Quality is the OPC UA status flattened to \code{GOOD}~(0),
\code{UNCERTAIN}~(1) and \code{BAD}~(2). Queries ignore bad samples unless
asked.

\subsection{Resampling: last value carried forward}

Samples arrive by exception, so between two stored samples the signal did not
drift; it was \emph{held}. Linear interpolation would invent a ramp that never
existed, and on a valve position or an alarm bit that is simply wrong. The
historian builds a regular grid and uses DuckDB's \code{ASOF LEFT JOIN} to
take, for every grid instant, the last sample at or before it:
\begin{equation}
  \hat x(t_g) = x(t_j), \qquad j = \max\{\,i : t_i \le t_g\,\}.
\end{equation}
Samples stored before the window's start count too, so a window never opens
with a column of nulls just because a tag happened to be quiet.

\begin{figure}[htbp]
\centering
\begin{tikzpicture}
\begin{axis}[width=0.6\textwidth,height=3.5cm,xlabel={time (s)},ylabel={value},
  xmin=0,xmax=10,ymin=0,ymax=5,grid=major,grid style={linegrey},legend pos=north west]
\addplot[only marks,mark=*,accent,mark size=2.5pt] coordinates {(0.4,1)(3.2,3)(6.7,2)(8.1,4)};
\addplot[slate,thick,const plot] coordinates {(0.4,1)(3.2,1)(3.2,3)(6.7,3)(6.7,2)(8.1,2)(8.1,4)(10,4)};
\addplot[badred,dashed] coordinates {(0.4,1)(3.2,3)(6.7,2)(8.1,4)};
\legend{stored samples, LOCF (used), interpolation (wrong)}
\end{axis}
\end{tikzpicture}
\caption{Step-hold resampling against linear interpolation.}
\label{fig:locf}
\end{figure}

\subsection{Covered span}
\code{span()} is the extent of the whole database, and a window opened there
almost always starts with nulls: tags report a few milliseconds apart, so the
earliest sample belongs to exactly one of them. \code{covered_span(tags)}
returns the interval over which \emph{every} requested tag has a value, the
latest of the per-tag first samples to the latest last sample. Use it to ask
``give me everything you have''.

\begin{warn}[Wall-clock timestamps]
Timestamps are wall clock, as a real historian's are. At \code{speed} $>1$ the
recorded timeline is compressed, so a window 60 wall seconds long covers
$60\times\text{speed}$ seconds of plant time. Training data comes from the
fast runner in plant time and never touches this database.
\end{warn}

\subsection{One writer}
DuckDB takes an exclusive lock on its file, so a second process cannot open
it even to read. That is not a limitation to work around: a historian is a
service you query, not a file you share. Exactly one process writes, and
everything else asks the query API on port 5404 (Chapter~\ref{ch:api}).

\section{The operator HMI}

The HMI (\code{pyplant/hmi/app.py}; port 5405 under the workbench, 8080
when started by hand) is a FastAPI application
serving one HTML page and a JSON API. It connects to the PLC over OPC UA,
polls the faceplate every 0.5\,s, feeds the alarm manager and reads trends
from the historian API. It shows a live mimic of the two tanks, the LIC101
faceplate (SP, PV, OP, mode), pump control, trends and the alarm list.

\subsection{Write guards}
Every operator write goes through one function that rejects bad requests
before anything reaches the control system:

\begin{center}
\begin{tabular}{@{}lll@{}}
\toprule
Check & HTTP status & Example \\
\midrule
tag exists & 404 & unknown tag \\
tag is operator-writable & 403 & writing \tg{LT101} \\
analog value inside span & 422 & setpoint of 900\,mm \\
connected to the PLC & 503 & PLC is down \\
\bottomrule
\end{tabular}
\end{center}

\subsection{Pulsed reset}
\code{POST /api/reset} pulses \tg{CMD_ALM_RESET}. A latched interlock needs
an \emph{edge}, not a level: holding the reset bit true would re-clear the
latch on every scan and defeat it.

\section{Alarm management (ISA-18.2)}

A boolean that turns a row red is not alarm management. ANSI/ISA-18.2 exists
because operators need to distinguish four situations that a single ``is it
on'' flag collapses into two. pyplant implements the core state model in
\code{pyplant/hmi/alarms.py}.

\begin{figure}[htbp]
\centering
\scalebox{0.8}{%
\begin{tikzpicture}[node distance=24mm and 44mm,
  st/.style={box,minimum width=34mm,minimum height=12mm}]
  \node[good,st] (n) {\textbf{NORMAL}\\inactive, acknowledged};
  \node[bad,st,right=of n] (u) {\textbf{UNACK\_ALARM}\\active, unacknowledged};
  \node[box,st,fill=warnamber!15,below=of u] (a) {\textbf{ACK\_ALARM}\\active, acknowledged};
  \node[box,st,fill=warnamber!25,below=of n] (r) {\textbf{RTN\_UNACK}\\cleared, unacknowledged};
  \node[box,st,fill=slate!15,above=12mm of $(n)!0.5!(u)$] (s) {\textbf{SHELVED}\\suppressed, expires};
  \draw[flow] (n) -- node[lbl,above]{condition in} (u);
  \draw[flow] (u) -- node[lbl,right]{acknowledge} (a);
  \draw[flow] (r) -- node[lbl,left]{acknowledge} (n);
  \draw[flow] ([xshift=-10mm]u.south) -- node[lbl,sloped,above=1pt,pos=0.35]{condition out} ([xshift=6mm]r.north);
  \draw[flow] ([xshift=16mm]r.north) -- node[lbl,sloped,below=1pt,pos=0.65]{in again} ([xshift=0mm]u.south);
  \draw[flow] (a.south) -- ++(0,-8mm) -| node[lbl,pos=0.25,below]{condition out} ([xshift=-8mm]r.west) |- (n.west);
  \draw[dflow] (s.west) -| node[lbl,pos=0.75,left]{unshelve / expiry} (n.north);
  \draw[dflow] (u.north) |- node[lbl,pos=0.25,right]{shelve} (s.east);
\end{tikzpicture}}
\caption{The ISA-18.2 alarm states implemented by \code{Alarm}. Shelving can be
entered from any state and returns to whatever the condition says at that
moment.}
\label{fig:isa}
\end{figure}

\begin{description}
  \item[NORMAL] quiet, and nobody needs to do anything.
  \item[UNACK\_ALARM] active, and nobody has looked at it yet.
  \item[ACK\_ALARM] active, and somebody has taken responsibility.
  \item[RTN\_UNACK] it cleared itself before anyone looked. This is the
        state that matters most: a transient between two operator glances
        leaves no trace in a flag-based list, and the transient nobody saw is
        often the interesting one. It stays on the list until acknowledged.
  \item[SHELVED] a deliberate, temporary, expiring silence with a reason
        recorded (default 30\,min, capped at 8\,h). The condition is still
        tracked underneath: if it goes active while shelved, the alarm comes
        back as unacknowledged on expiry.
\end{description}

\subsection{Priorities}
Priority is a design decision per alarm, not something derived from a name.
An interlock that trips the pump outranks a level approaching its limit.

\begin{center}
\begin{tabular}{@{}lll@{}}
\toprule
Alarm & Priority & Message \\
\midrule
\tg{ILK_PMP101} & CRITICAL (4) & interlock tripped, pump stopped, valve shut \\
\tg{ALM_LT101_LOLO} & HIGH (3) & tank 1 level low low \\
\tg{ALM_PT101_HIHI} & HIGH (3) & pump discharge pressure high high \\
\tg{ALM_LT101_HI} & MEDIUM (2) & tank 1 level high \\
\tg{ALM_HEALTH} & LOW (1) & model health index degraded \\
\bottomrule
\end{tabular}
\end{center}

The alarm page sorts by priority, then newest first. The summary reports the
counts per state, the number needing attention (UNACK plus RTN\_UNACK) and
the worst priority among them, which drives the banner colour.

\subsection{Operating the alarm list}
\begin{enumerate}
  \item A red row is an unacknowledged active alarm. Investigate, then press
        \ui{Ack}; it turns amber (acknowledged) and stays until the condition
        clears.
  \item A row marked \emph{returned} cleared on its own. Read it and
        acknowledge it; it then leaves the list.
  \item A chattering nuisance alarm can be shelved with a reason and a
        duration. Shelving is how you avoid operators ignoring the whole
        list instead.
  \item The interlock additionally needs \ui{Reset} once the cause is gone;
        acknowledging only clears the notification, not the trip.
\end{enumerate}
